\chapter{Технологическая часть}

\section{Средства реализации}
В качестве языка программирования был выбран Swift версии~6.4, так как он является языком со статической типизацией и основным языком разработки для операционной системы macOS, на которой выполнялась лабораторная работа. Сборка выполнялась с помощью Swift Package Manager.

При реализации использовались следующие средства:
\begin{itemize}
	\item тип \texttt{Array} стандартной библиотеки Swift~--- матрица представлена как массив строк \texttt{[[Int]]}, вспомогательные массивы $mulH$ и $mulV$~--- как \texttt{[Int]};
	\item функция \texttt{clock\_gettime\_nsec\_np} системной библиотеки macOS~\cite{clock-gettime}~--- для замера процессорного времени (подробнее в разделе~\ref{sec:time-tools});
	\item библиотека Swift Testing~--- для модульного тестирования реализаций.
\end{itemize}

\section{Реализация алгоритмов}
Тип матрицы и функция создания нулевой матрицы приведены в листинге~\ref{lst:matrix}, реализации алгоритмов~--- в листингах~\ref{lst:standard}--\ref{lst:winograd-shift-add}. Фрагменты кода, которые выводят промежуточные данные только в отладочной сборке (блоки \texttt{\#if DEBUG}), в листингах опущены.

% листинги готовит gen-listings.py: части в minipage не разрываются между страницами
% сгенерировано gen-listings.py
\noindent\begin{minipage}{\linewidth}
\lstinputlisting[caption={Тип матрицы и создание нулевой матрицы}, label={lst:matrix}]{matrix.swift}
\end{minipage}


% сгенерировано gen-listings.py
\noindent\begin{minipage}{\linewidth}
\lstinputlisting[caption={Стандартный алгоритм умножения матриц}, label={lst:standard}]{standard.swift}
\end{minipage}


% сгенерировано gen-listings.py
\noindent\begin{minipage}{\linewidth}
\lstinputlisting[caption={Алгоритм Винограда}, label={lst:winograd}]{winograd-1.swift}
\end{minipage}

\noindent\begin{minipage}{\linewidth}
\lstinputlisting[title={Продолжение листинга~\ref{lst:winograd}}]{winograd-2.swift}
\end{minipage}


% сгенерировано gen-listings.py
\noindent\begin{minipage}{\linewidth}
\lstinputlisting[caption={Алгоритм Винограда с оптимизацией А}, label={lst:winograd-shift}]{winograd-shift-1.swift}
\end{minipage}

\noindent\begin{minipage}{\linewidth}
\lstinputlisting[title={Продолжение листинга~\ref{lst:winograd-shift}}]{winograd-shift-2.swift}
\end{minipage}


% сгенерировано gen-listings.py
\noindent\begin{minipage}{\linewidth}
\lstinputlisting[caption={Алгоритм Винограда с оптимизацией Б}, label={lst:winograd-add}]{winograd-add-1.swift}
\end{minipage}

\noindent\begin{minipage}{\linewidth}
\lstinputlisting[title={Продолжение листинга~\ref{lst:winograd-add}}]{winograd-add-2.swift}
\end{minipage}


% сгенерировано gen-listings.py
\noindent\begin{minipage}{\linewidth}
\lstinputlisting[caption={Алгоритм Винограда с оптимизациями А и Б}, label={lst:winograd-shift-add}]{winograd-shift-add-1.swift}
\end{minipage}

\noindent\begin{minipage}{\linewidth}
\lstinputlisting[title={Продолжение листинга~\ref{lst:winograd-shift-add}}]{winograd-shift-add-2.swift}
\end{minipage}


\section{Тестирование}
Тестирование выполнялось методом чёрного ящика. Набор функциональных тестов, применявшийся ко всем пяти реализациям, приведён в таблице~\ref{tbl:tests}. Тесты 1--10 проверяют результат на заранее вычисленных значениях, тест 11 сравнивает результат каждой реализации с результатом стандартного алгоритма на случайных матрицах с элементами из отрезка $[-50; 50]$. Листинг модульных тестов приведён в приложении~А.
\begin{table}[H]
\caption{Функциональные тесты}
\label{tbl:tests}
\centering
\small
\setlength{\arraycolsep}{2.5pt}
\begin{tblr}{
  width = \linewidth,
  colspec = {Q[c,m] X[1.5,l,m] X[1.7,c,m] X[1.7,c,m] X[1.7,c,m] Q[c,m]},
  colsep = 3pt,
  hlines, vlines,
  row{1} = {font=\bfseries},
}
№ & Описание теста & $A$ & $B$ & Ожидаемый результат & Получ. \\
1 & Матрицы $1 \times 1$ & $(3)$ & $(-4)$ & $(-12)$ & Совп. \\
2 & Умножение на единичную матрицу $E$ слева и справа &
  $\begin{pmatrix}1&2&3\\4&5&6\\7&8&9\end{pmatrix}$ & $E_{3 \times 3}$ &
  $AE = EA = A$ & Совп. \\
3 & Умножение на нулевую матрицу &
  $\begin{pmatrix}1&2\\3&4\end{pmatrix}$ & $\begin{pmatrix}0&0\\0&0\end{pmatrix}$ &
  $\begin{pmatrix}0&0\\0&0\end{pmatrix}$ & Совп. \\
4 & Квадратные $2 \times 2$, чётное $m$ &
  $\begin{pmatrix}1&2\\3&4\end{pmatrix}$ & $\begin{pmatrix}5&6\\7&8\end{pmatrix}$ &
  $\begin{pmatrix}19&22\\43&50\end{pmatrix}$ & Совп. \\
5 & Квадратные $3 \times 3$, нечётное $m$ &
  $\begin{pmatrix}1&2&3\\4&5&6\\7&8&9\end{pmatrix}$ & $\begin{pmatrix}9&8&7\\6&5&4\\3&2&1\end{pmatrix}$ &
  $\begin{pmatrix}30&24&18\\84&69&54\\138&114&90\end{pmatrix}$ & Совп. \\
6 & Прямоугольные $2 \times 3$ и $3 \times 2$, нечётное $m$ &
  $\begin{pmatrix}1&2&3\\4&5&6\end{pmatrix}$ & $\begin{pmatrix}7&8\\9&10\\11&12\end{pmatrix}$ &
  $\begin{pmatrix}58&64\\139&154\end{pmatrix}$ & Совп. \\
7 & Прямоугольные $3 \times 2$ и $2 \times 4$, чётное $m$ &
  $\begin{pmatrix}1&2\\3&4\\5&6\end{pmatrix}$ & $\begin{pmatrix}1&2&3&4\\5&6&7&8\end{pmatrix}$ &
  $\begin{pmatrix}11&14&17&20\\23&30&37&44\\35&46&57&68\end{pmatrix}$ & Совп. \\
8 & Строка на столбец &
  $\begin{pmatrix}1&2&3&4\end{pmatrix}$ & $\begin{pmatrix}1&2&3&4\end{pmatrix}^T$ &
  $(30)$ & Совп. \\
9 & Столбец на строку, $m = 1$ &
  $\begin{pmatrix}1&2&3\end{pmatrix}^T$ & $\begin{pmatrix}4&5\end{pmatrix}$ &
  $\begin{pmatrix}4&5\\8&10\\12&15\end{pmatrix}$ & Совп. \\
10 & Отрицательные элементы &
  $\begin{pmatrix}-1&2\\3&-4\end{pmatrix}$ & $\begin{pmatrix}-5&-6\\7&8\end{pmatrix}$ &
  $\begin{pmatrix}19&22\\-43&-50\end{pmatrix}$ & Совп. \\
11 & Случайные матрицы, $m = \overline{1, 6}$, $n, q \in [1; 5]$ &
  Случайная $n \times m$ & Случайная $m \times q$ &
  Результат стандартного алгоритма & Совп. \\
\end{tblr}
\end{table}

Все тесты пройдены успешно.
